\documentclass[UTF8,12pt]{ctexart}     %点明文档类型是ctexart
\usepackage{indentfirst}
\usepackage{microtype}
\usepackage{float} 
\usepackage{amsmath}
\usepackage{booktabs}

\setlength{\parindent}{2em} % 设置首行缩进
\newtheorem{thm}{定理}
\numberwithin{thm}{section}  


\title{杂谈勾股定理}
\author{张三}
\date{\today}   % \today表示当前日期

\bibliographystyle{plain} % 参考文献的格式

\begin{document}    % begin和end标识出正文的范围

\maketitle   % 生成标题
\begin{abstract}
    这是一个关于勾股定理的文档
\end{abstract}
\tableofcontents   % 生成目录
\section{勾股定理在古代}

西方称勾股定理为毕达哥拉斯定理，将勾股定理的发现归功于公元前 6 世纪的
毕达哥拉斯学派。该学派得到了一个法则，可以求出可排成直角三角形的三元
数组。毕达哥拉斯学派没有书面著作，该定理的严格表述和证明则见于欧几里得\footnote{欧几里得，约公元前 330--275 年。}
《几何原本》的命题 47：“直角三角形斜边上的正方形等于两直角边上的两个正
方形之和。”证明是用面积做的。

我国《周髀算经》载商高（约公元前 12 世纪）答周公问：
\begin{quote}
    \zihao{-5}\kaishu
    勾广三，股修四，径隅五。
\end{quote}

又载陈子（约公元前 7--6 世纪）答荣方问：
\begin{quote}
    \zihao{-5}\kaishu
    若求邪至日者，以日下为勾，日高为股，勾股各自乘，并而开方除之，得邪至日。
\end{quote}

都较古希腊更早。后者已经明确道出勾股定理的一般形式。图 1 是我国古代对
勾股定理的一种证明。

\section{勾股定理的近代形式}
勾股定理可以用现代语言表述如下：
\begin{thm}[勾股定理]
    直角三角形斜边的平方等于两腰的平方和。

    可以用符号语言表述为:设直角三角形$ABC$,其中$\angle C = 90^\circ$,则有

    \[
        AB^2 = AC^2 + BC^2
    \]

\end{thm}

满足式 (1) 的整数称为\emph{勾股数}。
第 1 节所说的毕达哥拉斯学派得到的
三元数组就是勾股数。下表列出一些较小的勾股数：
\begin{table}[H]
    \begin{tabular}{|r|r|r|}
        \hline
        直角边 $a$ & 直角边 $b$ & 斜边 $c$ \\
        \hline
        3       & 4       & 5      \\
        5       & 12      & 13     \\
        \hline
    \end{tabular}
    \qquad
    ($a^2 + b^2 = c^2$)
\end{table}
\bibliography{math} % 参考文献的文件名

\end{document}